\noindent This appendix summarizes collection, preprocessing, and evaluation details for the datasets used in the paper.
For concrete examples of raw trajectory formats (asciinema \texttt{.cast} and \texttt{vhs} scripts), see Appendix~\ref{appendix:cligen-samples}. 
% 
Data collection uses a three-stage pipeline to maintain synchronized timing, privacy controls, and consistent, well-documented artifacts across CLIGen and GUIWorld.

\noindent\textbf{Sourcing.}
We construct our datasets from three complementary sources---public terminal recordings, scripted terminal replays, and a controlled desktop-capture rig.
Across all sources, we emphasize rights-respecting acquisition, privacy filtering, and temporally aligned interface signals (frames, actions, and text when available):
\begin{tcolorbox}[
  colback=metablue!8,
  colframe=metablue!40,
  boxrule=0pt,
  borderline west={1pt}{0pt}{gray!60},
  left=6pt,right=4pt,top=3pt,bottom=3pt]
\small
\begin{itemize}
  \setlength{\itemsep}{2pt}
  \setlength{\topsep}{2pt}
  \item \cligenGeneralLogo{}\,\textbf{CLIGen (General)}: public \texttt{asciinema} \texttt{.cast} archives; traces are replayed with official tools to preserve recorded terminal appearance (color schemes, cursor visibility, window geometry).
  \item \cligenCleanLogo{}\,\textbf{CLIGen (Clean)}: deterministic \texttt{vhs} scripts (e.g., package installs, REPLs, log filters) executed in isolated environments.
  \item \guiworldLogo{}\,\textbf{GUIWorld}: footage captured with the rig in \Cref{section:impl-guiworld}, pairing RGB video with low-latency pointer/key logs and optional accessibility cues (logged for analysis; not used as model inputs).
\end{itemize}
\end{tcolorbox}

\noindent\textbf{Alignment and sanitization.} 
All modalities share a common clock.
We align pointer/key events to the nearest frame, apply drift correction when needed, and drop clips with residual misalignment.
Privacy filters remove terminal sessions with sensitive strings and redact GUI regions likely to contain private content.
Frozen or repeated-frame recordings (capture artifacts) are discarded.

\noindent\textbf{Episode packaging.}
Runs are windowed into fixed-length, overlapping episodes (window sizes and strides are specified in the released configs).
Each shard stores RGB frames, terminal buffers or GUI metadata, serialized actions, the source tool (\texttt{asciinema}/\texttt{vhs}/GUI capture), and environment metadata.
Structured fields (buffers/metadata) are used for alignment and evaluation, but are not provided to the video models as state inputs.
Downstream dataloaders reconstruct batches directly.
Released configs specify preprocessing and windowing so external users can rebuild the corpus.

\subsection{Caption fields and metadata (CLIGen General)}
For CLIGen (General), each replayed \texttt{.cast} fragment (example shown in \Cref{appendix:cast-format}) is paired with three aligned descriptions and a compact metadata record.
Table~\ref{tab:cast-captions} summarizes the fields for clip \texttt{7****\_0001}.

\begin{smallschemetable}{Caption tiers and metadata for CLIGen (General) clip \texttt{7****\_0001}.}{tab:cast-captions}
  \rowcolors{2}{metablue!8}{white}
  \begin{tabularx}{\linewidth}{@{}l X@{}}
    \toprule
    \rowcolor{metablue!16}\textbf{Field} & \textbf{Content (abridged)} \\
    \texttt{caption} &
      The user (username field shown verbatim as \texttt{annonomous} [sic]) logged into their account on the host \texttt{NeuralComputer}, and after some initial terminal setup and cursor movements,
      they started typing the command \texttt{nvim} to likely open the Neovim text editor, but the session recording ends abruptly without showing the actual command execution or any further interaction. \\
    \texttt{caption\_detailed} &
      At the \texttt{annonomous@NeuralComputer:~\$} prompt, the user types \texttt{nvim}, but before the command executes, the screen rapidly displays a colorful,
      pixelated animation with various RGB colors, including shades of blue, green, and purple, filling the 80x24 terminal window.
      The animation is briefly interrupted by a cursor blink, then the prompt returns, awaiting further input. \\
    \texttt{caption\_semantic} &
      The terminal displays a rapid sequence of colored text and pixel art animations, with syntax-highlighted code and prompts appearing in a mix of blue, green,
      and white hues, as the user types commands and the system responds with scrolling output, error messages, and success indicators. \\
    \midrule
    \texttt{data\_info.version} & \texttt{2} (\texttt{asciinema} v2 header) \\
    \texttt{data\_info.size} & \texttt{width=80}, \texttt{height=24}; original/target/scaled sizes all match, padding \texttt{[0,0]}. \\
    \texttt{data\_info.env} & \texttt{SHELL=/bin/bash}, \texttt{TERM=xterm-256color}. \\
    \texttt{meta.videogen} &
      Automatically derived stats: visual complexity \texttt{2897.7}, interaction density \texttt{1.0}, color usage \texttt{518}, screen clears \texttt{6},
      event rate \texttt{7.8} events/s, typing rhythm (avg interval \texttt{0.13}s, variance \texttt{0.34}), and \texttt{38} visual-change events. \\
    \texttt{metadata} &
      Source recording on Asciinema: id \texttt{7****}, title \texttt{``aneo.nvim demo''}, author \texttt{annonomous} [sic], creation time \texttt{2025-05-13T23:37:53Z}, URLs for the page and raw \texttt{.cast}. \\
    \bottomrule
  \end{tabularx}
\end{smallschemetable}

These fields make each CLIGen (General) clip self-contained.
The three caption tiers provide prompts at different levels of detail.
\texttt{data\_info} and \texttt{metadata} preserve the structure needed to rebuild terminal geometry, environment, and source from the raw \texttt{.cast}. 
% A heuristic category count indicates that the dataset covers diverse terminal usage.
% It includes general terminal operations (610,331 clips), file operations (81,783), programming activities (66,430), system administration (42,668), and text editing (22,777).
% This distribution reflects typical command-line workflows across different user tasks.

\subsection{OCR evaluation protocol (CLIGen Clean)}
For CLIGen (Clean), OCR-based metrics evaluate how closely generated terminal videos match reference renderings derived from the ground-truth buffers in text space rather than pixels.
Each sample consists of a generated video and its paired reference video (matched by clip ID).
We keep only IDs where both videos are present. 
% 
From each paired video we use at most $K{=}5$ frames.
Let $T_\text{gen}$ and $T_\text{gt}$ be the frame counts of the generated and reference videos.
We set $T=\min(T_\text{gen},T_\text{gt})$.
If $T \le K$, we use all indices in $[0, T{-}1]$; otherwise, we select $K$ evenly spaced indices by deterministic rounding, then deduplicate and sort them so evaluation frames are spread across the trajectory.
For every selected index we read the corresponding frame from both videos.

Each frame is converted to RGB and passed to Tesseract OCR.
The resulting string is split into lines, leading and trailing whitespace is stripped, and internal whitespace is normalized by collapsing runs of spaces.
Empty lines are dropped.
We keep case and punctuation intact so that commands, paths, and symbols remain visible.
This gives an ordered list of normalized lines for the ground-truth frame ($g_1,\dots,g_{N_g}$) and the generated frame ($p_1,\dots,p_{N_p}$).

We summarize the OCR text-space metrics used per sampled frame as follows:
\begin{tcolorbox}[
  colback=metablue!8,
  colframe=metablue!40,
  boxrule=0pt,
  borderline west={1pt}{0pt}{gray!60},
  left=6pt,right=4pt,top=3pt,bottom=3pt]
\small
\textbf{Character accuracy.} This metric pools all lines into a single multi-line string for each side and measures normalized edit distance.
Let $s$ and $t$ be the concatenated ground-truth and generated texts and $d(s,t)$ their Levenshtein distance (insert/delete/replace cost $1$).
If both $s$ and $t$ are empty we set $\text{char\_acc}=1$; if only $s$ is empty we set $\text{char\_acc}=0$.
Otherwise,
\[
  \text{char\_acc} = \max\Bigl(0,\ 1 - \frac{d(s,t)}{\max(|s|,1)}\Bigr),
\]
Extra or missing characters are normalized by the reference length through the denominator $\max(|s|,1)$.
Frame-level scores are averaged over the selected frame pairs (up to $K{=}5$) to yield a per-video character accuracy, and group-level scores report the mean over videos.

\textbf{Exact-line accuracy.} This metric treats lines as position-sensitive units and reports a recall over ground-truth lines.
For a given frame, we compare line $g_i$ to $p_i$ at the same index.
A line is counted as correct only if $i\le N_p$ and $p_i=g_i$; lines that appear in the wrong position do not count.
If both lists are empty we set $\text{exact\_line\_acc}=1$; if the ground-truth list is empty but the generated list is not, we set $\text{exact\_line\_acc}=0$.
Otherwise,
\[
  \text{exact\_line\_acc} = \frac{1}{N_g}\sum_{i=1}^{N_g} \mathbf{1}[i \le N_p \land p_i = g_i].
\]
As with character accuracy, frame scores are averaged over the $K$ sampled frames to obtain a per-video score and then averaged over videos for the reported aggregate.
\end{tcolorbox} 
% Unless otherwise noted, OCR results in the main text use 1,000 randomly sampled video pairs (five frame pairs per video).
Together, these two metrics stress both fine-grained text fidelity and line-ordered terminal state reconstruction.

\subsection{Evaluation metrics and protocol (GUIWorld)}
\label{appendix:gui-metrics}
We report both global video metrics and action-driven metrics that focus on post-interaction frames. 
We compute these metrics using our GUIWorld evaluation suite.
We summarize the GUIWorld protocol at a glance as follows:
\begin{tcolorbox}[
  colback=metablue!8,
  colframe=metablue!40,
  boxrule=0pt,
  borderline west={1pt}{0pt}{gray!60},
  left=6pt,right=4pt,top=3pt,bottom=3pt]
\small
\begin{itemize}
  \setlength{\itemsep}{2pt}
  \setlength{\topsep}{2pt}
  \item \textbf{Global metrics} (\FVD$_\text{all}$, \SSIM$_\text{all}$, \LPIPS$_\text{all}$): computed over paired generated/ground-truth videos after standardized decoding, subsampling, and resizing.
  \item \textbf{Action-driven metrics} (\SSIM$_{+15}$, \LPIPS$_{+15}$, \FVD$_{+15}$): computed on post-action windows to measure interface fidelity after interaction events.
\end{itemize}
\end{tcolorbox}

\noindent\textbf{Global \FVD$_\text{all}$/\SSIM$_\text{all}$/\LPIPS$_\text{all}$}
We decode paired generated/ground-truth videos into RGB frames with temporal subsampling and resizing (\texttt{fps}=3, \texttt{size}=256, and \texttt{max\_seconds}=5 by default).
\SSIM$_\text{all}$~is computed using \texttt{torchmetrics} on frame tensors normalized to $[0,1]$ and averaged over frames.
\LPIPS$_\text{all}$~uses the AlexNet backbone on frames normalized to $[-1,1]$ and is averaged over frames.
\FVD$_\text{all}$~is computed in an \texttt{r3d18} embedding space (\texttt{prelogits} by default).
We extract features from fixed-length clips (16 frames at 112$\times$112 after uniform subsampling/padding).
We compute the Fr\'echet distance between the generated and reference feature distributions.

\noindent\textbf{Action-driven metrics (\SSIM$_{+15}$, \LPIPS$_{+15}$, \FVD$_{+15}$)}
For each paired rollout, we load recorded action timestamps (from JSON/CSV logs), map each timestamp $\tau$ to a frame index $f=\mathrm{round}(\tau \cdot \mathrm{fps})$, and clamp to the valid frame range.
We skip the action frame itself (\texttt{action\_start\_offset}=1) and evaluate the next $k{=}15$ frames after each action.
Concretely, for each action frame index $f$, we form the post-action set $\{f+1,\dots,f+k\}$ and clip it to valid frame indices.
We then take the deduplicated union over actions in the clip and keep frame indices in chronological order.
Clips with zero logged actions, or with empty valid post-action windows after clipping, are excluded from $+15$ metrics.
For \SSIM$_{+15}$~/\LPIPS$_{+15}$, we compute per-clip means over selected frame pairs and then average across clips.
For action-driven \FVD$_{+15}$, we build an \emph{after-action clip} per video by concatenating the same selected post-action frames, then uniformly subsample/pad to 16 frames and compute the Fr\'echet distance in the same \texttt{r3d18} feature space.
